% Generated by tools/textbook/export_listings.js from tools/concepts/assessments.json: do not edit.
\begin{chapstatus}{Il filo del capitolo}
I numeri sono diventati posizioni, colori e una velocità: l'intero blocco sull'informazione, in un gioco che gira. Prossimo passo: la logica che permetterà al giocatore di guidarlo.
\par\medskip\textbf{Biglietto d'uscita} (da due a quattro minuti, da soli, su carta):
\begin{enumerate}[nosep,leftmargin=2em]
\item Verm1 dichiara title: .asciz "VERM1" e head\_x: .word 5. In quale sezione va ciascuno, e perché head\_x non deve stare in .rodata?
\item step\_q8 cresce di SPEED\_Q8 a ogni frame, e la testa si sposta di una cella ogni volta che arriva a 256. Con SPEED\_Q8 = 64, quanti frame tra due passi? Con 128? Prevedete, poi eseguite.
\end{enumerate}
\end{chapstatus}
